% !TEX program = xelatex
\documentclass[UTF8,a4paper,11pt]{ctexart}

\usepackage{amsmath}
\usepackage{amssymb}
\usepackage{bm}
\usepackage{graphicx}
\usepackage{booktabs}
\usepackage{hyperref}

\newcommand{\Ind}[1]{\mathbf{1}\!\left[#1\right]}

\begin{document}

\section{Introduction}

高反射场景重建是神经渲染与新视角合成中的长期挑战。与近似朗伯表面不同，镜面反射、折射、近场照明以及复杂材质会产生强烈的 view-dependent effects，使同一空间区域在不同观察方向下呈现显著不同的颜色、高光形状和反射内容。这类效应在真实物体采集、数字资产生成、可编辑三维内容生产、增强现实和机器人感知中都十分常见，因此高质量反射场景重建不仅要求合成图像在像素层面逼真，还要求反射外观与底层几何、可见性和材质线索保持可解释的关系。

近年来，NeRF-like 方法、3D Gaussian Splatting、2D Gaussian Splatting 及一系列 Gaussian-based renderer 显著提升了新视角合成的质量和效率。然而，在高反射区域，常规多视角 photometric loss 并不自动区分由真实几何支撑的反射变化与由单视角颜色拟合得到的局部残差。模型可以通过 view-dependent color query、emission radiance、局部特征、法线扰动或不透明度调整来解释某一训练视角中的 specular appearance，却不必保证相邻视角中同一表面区域的反射预测具有稳定的几何对应关系。这种歧义使高光容易被当作逐视角颜色现象，而不是被约束为由表面、视线和可见性共同决定的跨视角结构。

为缓解这一问题，已有工作引入了方向查询、反射方向编码、deferred shading、材质分解和 directional factorization 等机制。这些方法改善了模型表达视角相关颜色的方式，使反射分量能够显式依赖观察方向、表面法线、粗糙度或 learned appearance features。Ref-GS 等方向分解类方法进一步分析了 orientation-viewing ambiguity，并通过更结构化的方向表示降低颜色、法线和方向查询之间的耦合。然而，这类方法主要回答的是“如何表示视角相关外观”，而不是“同一表面在不同视角中的反射预测是否具有几何一致性”。当监督仍主要来自最终 RGB 重建时，反射分量仍可能在不同视角间缺少显式约束。

我们观察到，反射外观虽然天然随视角变化，但许多错误的高光拟合表现为缺少几何支撑的局部单视角残差。若一个可微 renderer 能够输出 specular、depth、alpha、normal 和 confidence 等 renderer intermediate buffers，则可以在相邻视角之间通过几何投影建立同一表面区域的候选对应，并在这些对应上约束反射预测的稳定性。需要强调的是，本文并不假设同一表面点在不同视角中的真实 specular radiance 必须相等；我们约束的是可见、几何一致且置信度较高的相邻视角对应区域中的反射预测稳定性。这一约束是面向优化的正则化目标，而不是严格的物理守恒定律。

基于这一观察，本文提出 RC 模块组，即一种可插拔反射一致性优化框架。该框架采用 renderer-buffer-driven optimization 的形式，将反射一致性约束建立在图像空间 intermediate buffers 上，而不是建立在某一种特定三维表示或某个特定材质网络上。具体而言，RC 模块组包含 renderer-buffer interface、cross-view reflection consistency、confidence-weighted reflective correspondence、scheduled training integration、可选 material smoothness，以及 reflection-specific evaluation protocol。通过这些模块，方法把单视角 specular fitting 转化为受几何对应、可见性和反射置信度共同调制的跨视角正则项。

RC 模块组的关键优势在于其表示无关性。原则上，只要底层方法能够提供最终渲染颜色、specular 或 view-dependent reflection 分量、alpha/visibility、depth、normal 以及 material 或 specular confidence，便可接入这一统一接口。因此，该框架可作为 plug-and-play optimization module 接入 3DGS、2DGS、GaussianShader、3DGS-DR、NeRF-like reflective reconstruction 和 deferred neural rendering 等不同 renderer。与复制某一方向编码、光照网络或特定 PBR 分支不同，RC 模块组只依赖语义等价的 image-space buffers 和相机几何，从而将方法贡献从具体表示中解耦出来。

本文将 RC 的收益谨慎定位为增强跨视角反射一致性，而不是保证所有图像质量指标、几何指标或材质分解指标同步提升。由于反射区域的几何对应、遮挡和可见性本身具有不确定性，提升 reflection consistency 可能与 full-image PSNR、SSIM 或 LPIPS 之间存在 trade-off。因此，本文同时强调 reflection-specific evaluation protocol：cross-view reflection consistency、reflective-region image quality 和 full-image quality 应分开定义、分开报告，并在结论中保持各自的适用边界。

\paragraph{Contributions.}
本文的主要贡献如下：
\begin{itemize}
    \item 本文提出一个脱离特定表示的 renderer-buffer-driven 反射一致性优化框架，将跨视角反射约束建立在 renderer intermediate buffers 上，使 RC 模块组可作为可插拔优化组件接入不同可微 renderer。
    \item 本文设计一种基于深度反投影、目标视角投影与可微采样的 cross-view reflection consistency loss，用于约束几何一致区域中的反射预测稳定性。
    \item 本文提出 confidence-weighted reflective correspondence，将 alpha/visibility、depth consistency、normal agreement 和 specular confidence 统一为可靠性权重，从而降低遮挡、错误深度、低可见性区域和非反射区域对正则项的干扰。
    \item 本文给出 scheduled training integration 与 reflection-specific evaluation protocol，使反射一致性模块能够在明确边界下评估其收益，并避免将反射一致性改善误写为全局图像质量或几何质量的同步改善。
\end{itemize}

\section{Method}

\subsection{Overview}

本文方法的目标是在不改变底层三维表示定义的前提下，为高反射场景重建引入跨视角反射一致性正则。给定多视角图像和相机参数，基础 renderer 可以是 Gaussian splatting、neural field、deferred shading renderer 或其他可微场景表示。RC 模块组只作用在 renderer output buffer 层：它读取每个视角的颜色、反射分量、可见性、深度、法线和置信度缓冲，并在这些缓冲之间建立跨视角约束。

整体流程包含三个步骤。首先，对于 source view 和 target view，分别渲染得到 renderer intermediate buffers。其次，利用 source depth 将 source 像素反投影到三维空间，再投影到 target view，并通过 differentiable sampling 采样 target buffers，从而建立候选 reflective correspondence。最后，利用 alpha/visibility、depth consistency、normal agreement 和 specular confidence 构造可靠性权重，对几何一致区域中的 specular prediction 施加 confidence-weighted reflection consistency loss。训练时，该正则项在 warm-up 之后以固定频率启用，以降低早期错误几何对应对优化的影响，并控制额外 target rendering 成本。

\subsection{Renderer-Buffer Interface}

设多视角训练集合为
\begin{equation}
\mathcal{D}=\{(I_i,\Pi_i)\}_{i=1}^{N},
\end{equation}
其中 \(I_i\) 是第 \(i\) 个视角的观测图像，\(\Pi_i=(K_i,T_i)\) 表示相机内参与外参。这里 \(K_i\) 为内参矩阵，\(T_i\) 为从世界坐标到相机坐标的刚体变换。给定任意可微 renderer \(R_\theta\)，其输出定义为
\begin{equation}
R_\theta(\Pi_i)\rightarrow \mathcal{B}_i,
\end{equation}
其中 \(\theta\) 表示底层可学习场景表示的参数，\(\mathcal{B}_i\) 表示第 \(i\) 个视角的 renderer intermediate buffers。

RC 模块组要求的最小 buffer contract 写为
\begin{equation}
\mathcal{B}_i=
\{\hat{I}_i,S_i,C_i,A_i,D_i,N_i^r,N_i^d\}.
\end{equation}
其中 \(\hat{I}_i\) 是最终渲染 RGB；\(S_i\) 是 specular 或 view-dependent reflection 分量；\(C_i\) 是 specular/material confidence；\(A_i\) 是 alpha 或 visibility；\(D_i\) 是 depth；\(N_i^r\) 是 renderer-predicted normal；\(N_i^d\) 是由 depth 推导得到的 normal。所有量均定义在图像域 \(\Omega_i\) 上。

该接口不要求 \(R_\theta\) 采用某一种特定 Gaussian 参数化，也不要求使用特定方向编码、特定光照网络或特定材质分解。对于显式分解反射分量的方法，\(S_i\) 可以直接取 specular branch 的输出；对于没有显式 specular branch 的方法，\(S_i\) 可以由 view-dependent residual 或 deferred reflection 分量给出。类似地，\(C_i\) 可以来自 roughness、learned reflective confidence、specular activation 或其他可反映反射区域可靠性的图像空间信号。该设计使 RC 模块组与底层表示解耦，只依赖语义等价的 buffers 和相机几何。

\subsection{Cross-view Reflective Correspondence}

给定 source view \(s\) 和 target view \(t\)，我们希望在几何可见且置信度较高的区域中比较两视角的反射预测。对 source 像素 \(u=(x,y)\in\Omega_s\)，记其齐次像素坐标为 \(\bar{u}=(x,y,1)^\top\)。利用 source depth \(D_s(u)\) 进行反投影，得到世界坐标中的三维点
\begin{equation}
\mathbf{x}_s(u)=
T_s^{-1}
\begin{bmatrix}
D_s(u)K_s^{-1}\bar{u}\\
1
\end{bmatrix}.
\end{equation}
该点随后被投影到 target view：
\begin{equation}
\begin{aligned}
\tilde{\mathbf{p}}_t(u) &= K_tT_t\mathbf{x}_s(u),\\
v(u) &= \pi(\tilde{\mathbf{p}}_t(u)),\\
\hat{z}_t(u) &= \tilde{\mathbf{p}}_t^z(u),
\end{aligned}
\end{equation}
其中 \(\pi(\cdot)\) 表示透视除法，\(v(u)\) 是 target 图像中的连续坐标，\(\hat{z}_t(u)\) 是由 source 几何投影得到的 target 深度。

由于 \(v(u)\) 通常不是整数像素坐标，我们使用双线性 differentiable sampling 从 target buffers 中读取对应值：
\begin{equation}
\tilde{S}_{t\leftarrow s}(u)=
\operatorname{sample}(S_t,v(u)).
\end{equation}
同理可得到 \(\tilde{A}_{t\leftarrow s}(u)\)、\(\tilde{D}_{t\leftarrow s}(u)\) 和 \(\tilde{N}_{t\leftarrow s}(u)\)。该 correspondence 的作用不是假设跨视角 specular radiance 在物理上完全相同，而是为相邻视角中几何一致的可见区域提供反射预测稳定性约束。由此，缺少稳定三维对应支撑的局部高光残差会受到抑制，而真实视角相关变化仍可通过 \(S_i\)、\(C_i\) 和可见性权重进行调制。

\subsection{Confidence-weighted Reflection Consistency}

跨视角投影可能受到遮挡、深度误差、低透明度区域、法线不一致以及非反射材质的影响。因此，本文不对所有 source 像素等权施加反射一致性，而是先定义一组可靠性条件。有效投影 mask 为
\begin{equation}
M_{\mathrm{proj}}(u)=
\Ind{v(u)\in\Omega_t}\,
\Ind{\hat{z}_t(u)>0}.
\end{equation}
其中 \(\Ind{\cdot}\) 表示指示函数。

为排除被遮挡或投影到错误表面的点，我们定义 depth consistency mask：
\begin{equation}
M_{\mathrm{depth}}(u)=
\Ind{
\frac{|\tilde{D}_{t\leftarrow s}(u)-\hat{z}_t(u)|}
{\max(\hat{z}_t(u),\epsilon)}
<\tau_d
}.
\end{equation}
其中 \(\tau_d\) 是相对深度容忍阈值，\(\epsilon\) 是数值稳定项。visibility mask 定义为
\begin{equation}
M_{\alpha}(u)=
\Ind{A_s(u)>\tau_\alpha}\,
\Ind{\tilde{A}_{t\leftarrow s}(u)>\tau_\alpha}.
\end{equation}
用于保证 source 和 target 中均存在可见表面。

除几何和可见性外，RC 还使用 normal agreement 和反射置信度对对应进行连续加权。source 视角中的法线一致性定义为
\begin{equation}
G_n(u)=
\max(0,\langle N_s^r(u),N_s^d(u)\rangle),
\end{equation}
反射置信度定义为
\begin{equation}
G_c(u)=C_s(u).
\end{equation}
当 renderer 没有显式 roughness 时，\(C_s\) 可以由 learned reflective confidence、specular branch activation、view-dependent residual magnitude 或高光强度估计得到。本文只要求 \(C_s\) 表示反射预测的可靠性或反射区域倾向，而不将其强行解释为物理粗糙度。

综合上述条件，二值 mask 和连续权重为
\begin{equation}
M(u)=M_{\mathrm{proj}}(u)
M_{\mathrm{depth}}(u)
M_{\alpha}(u),
\end{equation}
\begin{equation}
w(u)=
M(u)\,
A_s(u)\,
\tilde{A}_{t\leftarrow s}(u)\,
G_n(u)\,
G_c(u).
\end{equation}
该权重使 RC loss 主要作用在可见、几何一致、法线可靠且具有较高反射置信度的对应区域中。

反射残差定义为
\begin{equation}
r_{s\rightarrow t}(u)=
\left\|
\operatorname{sg}(S_s(u))-
\tilde{S}_{t\leftarrow s}(u)
\right\|_1,
\end{equation}
其中 \(\operatorname{sg}(\cdot)\) 表示 stop-gradient。该设计减少 source 和 target 反射预测相互追逐造成的训练不稳定，使梯度主要作用于被采样的 target 反射预测。对称形式 \(\mathcal{L}_{s\rightarrow t}+\mathcal{L}_{t\rightarrow s}\) 可作为扩展，但并不是 RC 模块组的必要条件。

最终的 pairwise reflection consistency loss 为
\begin{equation}
\mathcal{L}_{\mathrm{RC}}(s,t)=
\frac{\sum_{u\in\Omega_s}w(u)r_{s\rightarrow t}(u)}
{\sum_{u\in\Omega_s}w(u)+\epsilon}.
\end{equation}
在训练中，对满足视角间隔、重叠区域和可见性条件的 view-pair distribution \(\mathcal{P}\) 取期望：
\begin{equation}
\mathcal{L}_{\mathrm{RC}}=
\mathbb{E}_{(s,t)\sim\mathcal{P}}
\left[
\mathcal{L}_{\mathrm{RC}}(s,t)
\right].
\end{equation}
当一个 view pair 中没有足够可靠的对应时，分母中的 \(\epsilon\) 保证目标函数数值稳定，并自然降低该 pair 对优化的影响。

\subsection{Optimization Objective and Training Schedule}

RC 模块组作为附加正则项接入基础 renderer 的训练目标。总体优化目标写为
\begin{equation}
\mathcal{L}=
\mathcal{L}_{\mathrm{rgb}}
+
\lambda_{\mathrm{RC}}(k)\mathcal{L}_{\mathrm{RC}}
+
\lambda_{\mathrm{reg}}(k)\mathcal{L}_{\mathrm{reg}},
\end{equation}
其中 \(k\) 表示训练迭代步，\(\mathcal{L}_{\mathrm{rgb}}\) 是基础方法原有的图像重建项，\(\mathcal{L}_{\mathrm{reg}}\) 表示可选正则项。作为常见选择，图像重建项可写为
\begin{equation}
\mathcal{L}_{\mathrm{rgb}}
=
(1-\lambda_{\mathrm{ssim}})
\|\hat{I}-I\|_1
+
\lambda_{\mathrm{ssim}}
\mathcal{L}_{\mathrm{DSSIM}}(\hat{I},I),
\end{equation}
但 RC 模块组并不要求所有 renderer 必须采用这一具体形式。

由于训练早期的 depth、alpha 和 normal 往往尚不稳定，过早施加跨视角几何对应可能引入错误梯度。本文采用 scheduled RC training integration，在 warm-up 后以固定频率启用 RC：
\begin{equation}
\lambda_{\mathrm{RC}}(k)=
\begin{cases}
0, & k<k_0,\\
\lambda_{\mathrm{RC}}, & k\ge k_0 \ \mathrm{and}\ k\bmod K=0,\\
0, & \mathrm{otherwise}.
\end{cases}
\end{equation}
其中 \(k_0\) 是 RC 启动迭代，\(K\) 是启用间隔。若每 \(K\) 次迭代计算一次 RC，则平均额外 target render 成本约为 \(1/K\)。这一调度使基础颜色、可见性和几何缓冲先获得相对稳定的初始化，再逐步引入 cross-view reflection consistency。

\subsection{Auxiliary Regularization and Geometry Filtering}

除核心 RC loss 外，RC 模块组可以包含辅助的 material 或 roughness smoothness 正则。给定反射置信度或材质缓冲 \(C\)，可定义局部平滑项
\begin{equation}
\mathcal{L}_{\mathrm{mat}}=
\sum_{u}
\left(
|\nabla_x C(u)|+|\nabla_y C(u)|
\right).
\end{equation}
若 renderer 显式输出 roughness map \(R\)，也可使用
\begin{equation}
\mathcal{L}_{\mathrm{rough}}=\operatorname{TV}(R).
\end{equation}
该项用于鼓励材质或反射置信度在局部上更平滑，减少孤立噪声，并可作为消融控制。它并不是 cross-view reflection consistency 的核心来源，也不应被表述为 RC 模块组的主要理论贡献。

在需要深度融合或表面提取的场景中，RC buffers 还可用于 confidence-aware geometry filtering。我们定义几何可信度
\begin{equation}
Q(u)=A(u)\max(0,\langle N^r(u),N^d(u)\rangle).
\end{equation}
在深度融合或表面提取前，可使用 \(Q(u)\) 过滤低可见性或法线不一致区域，以减少不可靠深度对几何后处理的影响。该模块与 RC loss 共享 alpha 和 normal agreement 的可靠性思想，但它属于工程增强模块。除非有独立几何指标支持，不能由此推出 mesh quality 会随之改善。

\subsection{Reflection-specific Evaluation Protocol}

为了评估反射一致性优化框架，本文采用 reflection-specific evaluation protocol，将反射一致性、反射区域图像质量和全图图像质量分开定义。第一类指标是 cross-view reflection consistency：在一组 view pairs 上计算与 \(\mathcal{L}_{\mathrm{RC}}\) 同型的误差，数值越低表示几何一致区域中的预测反射分量越稳定。

第二类指标是 reflective-region image quality。给定由 \(C\)、alpha 或 specular intensity 定义的 reflective mask \(M_{\mathrm{refl}}\)，可在反射区域内报告 PSNR、SSIM 和 LPIPS 等指标：
\begin{equation}
M_{\mathrm{refl}}(u)=
\Ind{A(u)>\tau_\alpha}\,
\Ind{C(u)>\tau_c}.
\end{equation}
其中 \(\tau_c\) 是反射置信度阈值。该类指标关注反射区域的图像重建质量，但不与跨视角反射稳定性等价。

第三类指标是 full-image quality，即在全图范围内报告常规 PSNR、SSIM 和 LPIPS。三类指标必须分开报告，因为提升 reflection consistency 不必然等价于提升全图图像质量，也不必然推出几何或材质分解质量同步提升。因此，RC 模块组应被评价为反射一致性优化框架，而不是保证所有图像质量指标全面提升的通用插件。

\end{document}
